%============================================================
% CH1 · PATTERN 1 — EXERCISES, HINTS, ANSWERS
%============================================================
\subsection{Pattern 1: Exercises and Going Further}
\label{ssec:ch1-p1-ex}

\begin{patternbox}{Pattern 1 Summary Card}
\begin{tabularx}{\linewidth}{@{}>{\bfseries}l X@{}}
Recognise & one number (or one global number per position) that depends on the values, not on pairs far apart; ``$\bigO{n}$ time, $\bigO{1}$ memory'' \\
Core loop & state $\gets$ summary of the empty prefix; for each $x$: state $\gets$ \textsc{update}(state, $x$); read the answer \\
States & sum / count; min / max with first or last index; top-$k$ (distinct or not); prefix extreme as the partner of later elements; current and best streak \\
Invariant & before step $i$ the state is exactly the summary of $[0, i)$ \\
Cost & $\bigO{n}$ time ($\bigO{nk}$ for top-$k$), $\bigO{1}$ extra memory; two passes when a global value is needed at every position \\
Top traps & maximum started at $0$; forgetting to shift; a sentinel that is a legal value; distinct vs elements; using the prefix after adding the current element \\
\end{tabularx}
\end{patternbox}

\exlegend
\begin{enumerate}[leftmargin=*, label=\textbf{E1.\arabic*}, itemsep=6pt]\sloppy
    \item \exeasy\ \pclassic \exprob{Max Consecutive Ones}{LeetCode 485}{https://leetcode.com/problems/max-consecutive-ones/}\ \lc{Easy}. A binary array of length $n \le 10^5$ is given. Return the length of the longest block of consecutive 1's.
    \item \exeasy\ \pconcept \exprob{Equal Candies}{Codeforces 1676B}{https://codeforces.com/problemset/problem/1676/B}\ \cf{800}. There are $n \le 50$ boxes with $a_i \le 10^7$ candies. You may only \emph{eat} candies (remove them from boxes). Find the minimum number eaten so that all boxes end with the same number of candies. Up to $1000$ test cases.
    \item \exeasy\ \pconcept \exprob{Maximum Difference Between Increasing Elements}{LeetCode 2016}{https://leetcode.com/problems/maximum-difference-between-increasing-elements/}\ \lc{Easy}. Return the maximum of $a_j - a_i$ over all $i < j$ with $a_i < a_j$, or $-1$ if no such pair exists. $2 \le n \le 1000$.
    \item \exeasy\ \precent \exprob{Guess the Maximum}{Codeforces 1979A}{https://codeforces.com/problemset/problem/1979/A}\ \cf{800}. Alice picks an integer $k$; then Bob chooses $i < j$ and takes $M = \max(a_i, \dots, a_j)$. Alice wins if $M > k$. Find the largest $k$ for which Alice wins whatever Bob does. $\sum n \le 5\cdot10^4$, $n \ge 2$.
    \item \exmedium\ \ptricky \exprob{Maximum Product of Three Numbers}{LeetCode 628}{https://leetcode.com/problems/maximum-product-of-three-numbers/}\ \lc{Easy}. Return the maximum product of three elements at distinct positions; $3 \le n \le 10^4$, $-1000 \le a_i \le 1000$. Do it in one pass.
    \item \exmedium\ \ptricky \exprob{Count Hills and Valleys in an Array}{LeetCode 2210}{https://leetcode.com/problems/count-hills-and-valleys-in-an-array/}\ \lc{Easy}. Index $i$ is part of a hill (valley) if its closest \emph{non-equal} neighbours on both sides exist and are both smaller (larger). Adjacent equal indices belong to the same hill or valley. Count the hills and valleys.
    \item \exmedium\ \pclassic \exprob{Removing Minimum and Maximum From Array}{LeetCode 2091}{https://leetcode.com/problems/removing-minimum-and-maximum-from-array/}\ \lc{Medium}. The values are distinct. One deletion removes the first or the last element. Return the minimum number of deletions that remove both the minimum and the maximum. $n \le 10^5$.
\end{enumerate}

\subsubsection*{Hints}
\begin{enumerate}[leftmargin=*, label=\textbf{E1.\arabic*}]
    \item Two variables: the length of the block ending at the current element, and the best length so far. What happens to the first one at a 0?
    \item Every box must end at the same level $L$, and you cannot add candies. What is the largest possible $L$? Then the answer is a sum.
    \item Fix the larger element $a_j$. Its best partner is an extreme of the prefix $[0, j)$. Mind the strict inequality.
    \item Bob wants $M$ as small as possible. Does he ever profit from a segment longer than two elements?
    \item Sorted thinking: the best product uses the three largest values, \emph{or} the largest value and the two most negative. Keep five slots.
    \item Equal neighbours are one plateau. Keep the last value that differed from the current one, and only look at an index when the next value differs.
    \item Let $i < j$ be the positions of the two extremes. Both can go from the front, both from the back, or one from each side.
\end{enumerate}

\subsubsection*{Answers}

\begin{solutionbox}[E1.1 Max Consecutive Ones]
State: \texttt{run} $=$ length of the block of 1's ending at the current position, \texttt{best} $=$ the longest block seen. A 1 extends the block, a 0 resets it.
\begin{lstlisting}[style=cpp]
#include <bits/stdc++.h>
using namespace std;

class Solution {
public:
    int findMaxConsecutiveOnes(vector<int>& nums) {
        int run = 0, best = 0;
        for (int x : nums) {
            run = (x == 1) ? run + 1 : 0;      // extend or reset the current block
            best = max(best, run);             // update the answer DURING the scan
        }
        return best;
    }
};

int main() {
    vector<vector<int>> tests = {{1, 1, 0, 1, 1, 1}, {1, 0, 1, 1, 0, 1}};
    for (auto& t : tests) cout << Solution().findMaxConsecutiveOnes(t) << " ";   // 3 2
    cout << "\n";
}
\end{lstlisting}
\begin{center}
\begin{tikzpicture}
    \dparray{o}{0.6}{nums}{1/awin, 1/awin, 0/aout, 1/akeep, 1/akeep, 1/akeep}
    \dparray{r}{-0.25}{\texttt{run}}{1/compcell, 2/compcell, 0/compcell, 1/compcell, 2/compcell, 3/anscell}
\end{tikzpicture}
\end{center}
The answer must be updated inside the loop: reading \texttt{run} only at the end returns the length of the \emph{last} block.
\end{solutionbox}

\begin{solutionbox}[E1.2 Equal Candies]
The final level $L$ cannot exceed any box, so $L \le \min a_i$; eating more than necessary only costs, so $L = \min a_i$. The answer is $\sum_i (a_i - L) = S - n \cdot \min a_i$: a sum and a minimum, one pass.
\begin{lstlisting}[style=cpp]
#include <bits/stdc++.h>
using namespace std;

int main() {
    int t;
    cin >> t;
    while (t--) {
        int n;
        cin >> n;
        long long sum = 0, mn = LLONG_MAX;
        for (int i = 0; i < n; i++) {
            long long x;
            cin >> x;
            sum += x;
            mn = min(mn, x);
        }
        cout << sum - n * mn << "\n";
    }
}
\end{lstlisting}
Official sample: $10, 4975, 38, 0, 0$ (for instance $1000\cdot5 + 5 - 6\cdot5 = 4975$). Here $50\cdot10^7$ still fits in \texttt{int}, but the habit costs nothing.
\end{solutionbox}

\begin{solutionbox}[E1.3 Maximum Difference Between Increasing Elements]
For each $j$, the best $i < j$ is the one with the smallest $a_i$: keep \texttt{lo} $= \min(a_0, \dots, a_{j-1})$. Use it \emph{before} adding $a_j$, and only if $a_j > \texttt{lo}$ (strictly), otherwise the pair is not ``increasing''.
\begin{lstlisting}[style=cpp]
#include <bits/stdc++.h>
using namespace std;

class Solution {
public:
    int maximumDifference(vector<int>& nums) {
        int lo = nums[0], best = -1;
        for (int j = 1; j < (int)nums.size(); j++) {
            if (nums[j] > lo) best = max(best, nums[j] - lo);   // pair (argmin of prefix, j)
            lo = min(lo, nums[j]);                              // now include nums[j]
        }
        return best;
    }
};

int main() {
    vector<vector<int>> tests = {{7, 1, 5, 4}, {9, 4, 3, 2}, {1, 5, 2, 10}};
    for (auto& t : tests) cout << Solution().maximumDifference(t) << " ";   // 4 -1 9
    cout << "\n";
}
\end{lstlisting}
On $[9, 4, 3, 2]$ every element is a new minimum, no pair qualifies, and the answer stays $-1$ --- a version that writes \texttt{best = max(best, nums[j] - lo)} unconditionally would return $-1$ only by luck (its best difference is $-1$ here, but it is $0$ on $[5,5]$ where the expected answer is $-1$). Checked against the $\bigO{n^2}$ double loop on 30000 random arrays.
\end{solutionbox}

\begin{solutionbox}[E1.4 Guess the Maximum]
Extending a segment can only increase its maximum, so Bob's best choices are segments of length two: $\min_i \max(a_i, a_{i+1})$. Alice wins exactly when $k$ is below this value, so the answer is $\min_i \max(a_i, a_{i+1}) - 1$ --- an aggregate (a minimum) over adjacent pairs.
\begin{lstlisting}[style=cpp]
#include <bits/stdc++.h>
using namespace std;

int main() {
    int t;
    cin >> t;
    while (t--) {
        int n;
        cin >> n;
        vector<int> a(n);
        for (int& x : a) cin >> x;
        int best = INT_MAX;
        for (int i = 0; i + 1 < n; i++) best = min(best, max(a[i], a[i + 1]));
        cout << best - 1 << "\n";
    }
}
\end{lstlisting}
Official sample: $3, 1, 0, 15, 9, 2$; for $(2, 4, 1, 7)$ the pair maxima are $4, 4, 7$, so the answer is $3$. The trap is to think about the global maximum, or to forget that Bob must take at least two elements ($i < j$). Checked against trying every segment for random arrays with $n \le 9$.
\end{solutionbox}

\begin{solutionbox}[E1.5 Maximum Product of Three Numbers]
If the array were sorted, the answer would be $\max(m_1 m_2 m_3,\ m_1 s_1 s_2)$, with $m_1 \ge m_2 \ge m_3$ the three largest and $s_1 \le s_2$ the two smallest elements: two negative factors make a positive product. The tempting ``three largest'' gives $6$ on $[-10, -9, 1, 2, 3]$ instead of $3\cdot(-10)\cdot(-9) = 270$. All five values come from one pass with two shifting lists (duplicates allowed, since positions must differ, not values).
\begin{lstlisting}[style=cpp]
#include <bits/stdc++.h>
using namespace std;

class Solution {
public:
    int maximumProduct(vector<int>& nums) {
        int m1 = INT_MIN, m2 = INT_MIN, m3 = INT_MIN;   // three largest
        int s1 = INT_MAX, s2 = INT_MAX;                 // two smallest
        for (int x : nums) {
            if (x > m1)      { m3 = m2; m2 = m1; m1 = x; }
            else if (x > m2) { m3 = m2; m2 = x; }
            else if (x > m3) m3 = x;
            if (x < s1)      { s2 = s1; s1 = x; }
            else if (x < s2) s2 = x;
        }
        return max(m1 * m2 * m3, m1 * s1 * s2);         // |product| <= 10^9 fits in int
    }
};

int main() {
    vector<vector<int>> tests = {{1, 2, -3}, {1, 2, 3, 4}, {-1, -2, -3}, {-10, -9, 1, 2, 3}};
    for (auto& t : tests) cout << Solution().maximumProduct(t) << " ";   // -6 24 -6 270
    cout << "\n";
}
\end{lstlisting}
Here \texttt{INT\_MIN}/\texttt{INT\_MAX} sentinels are safe: $n \ge 3$, so every slot is filled with a real value before it is used. Checked against all triples on 30000 random arrays.
\end{solutionbox}

\begin{solutionbox}[E1.6 Count Hills and Valleys in an Array]
Treat each plateau of equal values as one point. Keep \texttt{prev}, the last value that differed from the current plateau. At index $i$ (with $0 < i < n-1$), skip it if $a_i = a_{i+1}$ (the plateau continues and will be judged at its last index); otherwise $a_i$ is a hill if $a_i > \texttt{prev}$ and $a_i > a_{i+1}$, a valley if both are smaller, and then \texttt{prev} becomes $a_i$.
\begin{lstlisting}[style=cpp]
#include <bits/stdc++.h>
using namespace std;

class Solution {
public:
    int countHillValley(vector<int>& nums) {
        int n = nums.size(), cnt = 0, prev = nums[0];   // prev: last different value on the left
        for (int i = 1; i + 1 < n; i++) {
            if (nums[i] == nums[i + 1]) continue;       // inside a plateau: judge it at its end
            if ((nums[i] > prev && nums[i] > nums[i + 1]) ||
                (nums[i] < prev && nums[i] < nums[i + 1])) cnt++;
            prev = nums[i];
        }
        return cnt;
    }
};

int main() {
    vector<vector<int>> tests = {{2, 4, 1, 1, 6, 5}, {6, 6, 5, 5, 4, 1}};
    for (auto& t : tests) cout << Solution().countHillValley(t) << " ";   // 3 0
    cout << "\n";
}
\end{lstlisting}
\begin{center}
\begin{tikzpicture}
    \dparray{h}{0}{nums}{2/acell, 4/akeep, 1/aread, 1/aread, 6/akeep, 5/acell}
    \node[note, anchor=west, text width=4.4cm, align=left] at (5.8,0) {hill at 4, one valley for the plateau $1,1$, hill at 6: answer 3};
\end{tikzpicture}
\end{center}
The trap: comparing with the immediate neighbours counts the plateau $1,1$ zero times (each 1 has an equal neighbour) or twice. Checked against ``compress equal neighbours, then compare'' on 30000 random arrays with small values.
\end{solutionbox}

\begin{solutionbox}[E1.7 Removing Minimum and Maximum From Array]
One pass finds the positions of the minimum and the maximum; call them $i < j$ (swap if needed). Deleting from the front up to $j$ removes both ($j + 1$ deletions); deleting from the back down to $i$ removes both ($n - i$); or delete the front up to $i$ and the back down to $j$ ($(i + 1) + (n - j)$). The answer is the minimum of the three.
\begin{lstlisting}[style=cpp]
#include <bits/stdc++.h>
using namespace std;

class Solution {
public:
    int minimumDeletions(vector<int>& nums) {
        int n = nums.size(), lo = 0, hi = 0;           // positions of min and max
        for (int k = 1; k < n; k++) {
            if (nums[k] < nums[lo]) lo = k;
            if (nums[k] > nums[hi]) hi = k;
        }
        int i = min(lo, hi), j = max(lo, hi);
        return min({j + 1, n - i, (i + 1) + (n - j)});
    }
};

int main() {
    vector<vector<int>> tests = {{2, 10, 7, 5, 4, 1, 8, 6}, {0, -4, 19, 1, 8, -2, -3, 5}, {101}};
    for (auto& t : tests) cout << Solution().minimumDeletions(t) << " ";   // 5 3 1
    cout << "\n";
}
\end{lstlisting}
\begin{center}
\begin{tikzpicture}
    \dpindex{ri}{-0.05}{}{0}{7}
    \dparray{r}{0.6}{nums}{2/aout, 10/aout, 7/acell, 5/acell, 4/acell, 1/aout, 8/aout, 6/aout}
    \aptrup{r1}{max}{accentorange} \aptrup{r5}{min}{successgreen}
    \node[note, anchor=west, text width=3.8cm, align=left] at (7.5,0.6) {Example 1: $i = 1$, $j = 5$: front-only $6$, back-only $7$, both sides $2 + 3 = 5$};
\end{tikzpicture}
\end{center}
Checked against trying every pair (front count, back count) on random arrays of distinct values.
\end{solutionbox}
